\section{Introduction}

Peer review stands as the cornerstone of scientific advancement, serving as the primary mechanism through which the academic community validates research quality and determines the dissemination of knowledge. Despite its central importance, traditional peer review systems face mounting challenges: reviewer shortages, prolonged review cycles, inconsistent evaluation standards across disciplines, and the inherent subjectivity of human judgment \citep{mackenzie2021peer}. These challenges have spurred growing interest in AI-assisted review systems that can augment---rather than replace---human expertise.

Recent years have witnessed the emergence of various AI tools for academic paper evaluation, including citation-based matching systems \citep{kang2018peerread}, topic modeling approaches \citep{li2023reviewerfinder}, and large language model-based review assistants \citep{stanford2024ai}. However, a critical gap persists: these tools predominantly operate at the ``technical'' level (\textit{shu}), implementing operational metrics without grounding in the underlying ``philosophical principles'' (\textit{dao}). The evaluation dimensions we take for granted---novelty, methodological rigor, evidence confidence, and practical impact---are not arbitrary criteria but products of decades of development in the philosophy of science, statistics, and the sociology of science. Their theoretical origins remain largely unexplored in existing AI systems.

More fundamentally, existing AI-assisted review tools suffer from a \textit{static knowledge problem}. They deploy fixed evaluation criteria that cannot adapt to the evolving standards of academic disciplines or learn from the feedback provided by editors and reviewers. This limitation stands in stark contrast to the emerging paradigm of \textit{self-improving agents} in the broader AI community. Systems like Reflexion \citep{shinn2023reflexion} and Voyager \citep{wang2023voyager} have demonstrated that agents can learn from their interactions, accumulating experience and refining their strategies over time. Yet these self-improving frameworks have been applied primarily to verifiable tasks such as code generation and game playing, leaving the subjective domain of academic evaluation largely unexplored.

In this paper, we present WuYa, a theory-driven multi-agent system that addresses both gaps simultaneously. WuYa makes the following contributions:

\begin{enumerate}
    \item \textbf{Retrieval-Evaluation Coupling}: We introduce a mechanism where retrieval-augmented generation (\rag{}) serves not merely as knowledge supplementation but as the \textit{core triggering mechanism} for evaluation. When a \subagent{} needs to explain \textit{why} a paper's innovation is limited or \textit{how} to improve its methodology, it retrieves relevant passages from canonical texts, transforming abstract principles into actionable guidance.

    \item \textbf{Two-Phase Routing Architecture}: We design a two-phase routing system where a CUDOS (Communism, Universalism, Disinterestedness, Organized Skepticism) \subagent{} serves as a gatekeeper, ensuring compliance with fundamental norms of scientific conduct before expert evaluation proceeds in parallel.

    \item \textbf{Hybrid Knowledge Strategy}: We implement a dual-layer knowledge architecture: principles (``dao'') from classical theories are internalized in system prompts, while evidence and examples (``shu'') are retrieved on-demand through \rag{}, combining the consistency of internalized knowledge with the flexibility of dynamic retrieval.

    \item \textbf{LLM-as-Mapper with Discipline Priors}: We propose a novel paper localization approach where an LLM maps five-dimensional evaluation scores to estimated journal tier (Q1-Q4, A-D), calibrated by discipline-specific prior distributions derived from SpringRank journal rankings. This addresses the cross-disciplinary generalization challenge inherent in purely statistical approaches.

    \item \textbf{Self-Improving Frontier Discovery}: Inspired by Reflexion's verbal reinforcement learning and Voyager's skill library, we introduce a mechanism where the system learns evaluation preferences from editor feedback, dynamically updating a ``frontier surface library'' that captures what each journal values. This transforms evaluation from a static function to an adaptive, learning-based process.
\end{enumerate}

To the best of our knowledge, WuYa represents the first systematic integration of (1) retrieval-augmented evaluation, (2) theory-driven multi-agent architecture, and (3) self-improving feedback learning for academic paper review. We validate WuYa on a benchmark of 5,000 papers across computer science, biology, economics, physics, and medicine, demonstrating significant improvements over baselines in both recommendation accuracy and review quality.

The remainder of this paper is organized as follows. Section~\ref{sec:related} reviews related work on AI-assisted review, multi-agent systems, \rag{}, and self-improving agents. Section~\ref{sec:theory} establishes the theoretical foundations by tracing each evaluation dimension to its philosophical origins. Section~\ref{sec:architecture} details the system architecture, including the two-phase routing and \subagent{} designs. Section~\ref{sec:localization} presents the paper localization methods. Section~\ref{sec:frontier} introduces the self-improving frontier discovery mechanism. Section~\ref{sec:implementation} provides implementation details. Section~\ref{sec:experiments} reports experimental results. Section~\ref{sec:discussion} discusses implications, limitations, and future directions.
